% Auto-extracted from: lecture-7-handout.pdf
% Backend: claude
% Review formulas marked with % OCR_UNCERTAIN before proceeding.

\documentclass{article}
\usepackage{amsmath,amssymb,amsthm}
\newtheorem{theorem}{Theorem}
\newtheorem{lemma}[theorem]{Lemma}
\newtheorem{corollary}[theorem]{Corollary}
\newtheorem{proposition}[theorem]{Proposition}
\newtheorem{definition}[theorem]{Definition}
\begin{document}

\begin{definition}[Fisher Information]
Suppose $\mathcal{P} = \{f_\theta : \theta \in \Theta\}$ where $f_\theta(x)$ is a p.d.f. with parameter $\theta$ w.r.t. $\nu$ and $\Theta$ is an open subset of $\mathbb{R}$. Suppose for any $\theta \in \Theta$, $\frac{\partial f_\theta(x)}{\partial\theta}$ exists and is finite, $P_\theta$-a.s. Let $X$ be a sample from $P_\theta \in \mathcal{P}$. To measure the amount of information that an observation $X$ carries about $\theta$, we look at the Fisher information defined as
\[I(\theta) = \mathbb{E}\left(\frac{\partial}{\partial\theta}\log f_\theta(X)\right)^2 = \int\left(\frac{\partial}{\partial\theta}\log f_\theta(X)\right)^2 f_\theta(X)d\nu(x)\]
\end{definition}

\begin{theorem}[Poisson Families]
Suppose $(X_1,\ldots,X_n)$ is a i.i.d. sample from a Poisson distribution $\mathcal{P}(\lambda)$. Then the joint PMF is
\[f_\lambda(x) = \prod_{i=1}^n \frac{\lambda^{x_i}}{x_i!}\exp(-\lambda)\]

We have:
- $\log f_\lambda(x) = \sum_i x_i \log(\lambda) - n\lambda - \sum_i \log(x_i!)$
- $\frac{\partial}{\partial\lambda} \log f_\lambda(x) = \frac{\sum_{i=1}^n x_i}{\lambda} - n$
- $I(\lambda) = \operatorname{Var}\left(\frac{\sum_{i=1}^n X_i}{\lambda}\right) = \frac{n\lambda}{\lambda^2} = \frac{n}{\lambda}$

Here we've used $\mathbb{E}\frac{\partial}{\partial\lambda}\log f_\lambda(X) = 0$ and thus $I(\lambda) = \operatorname{Var}\left(\frac{\partial}{\partial\lambda}\log f_\lambda(X)\right)$
\end{theorem}

\begin{theorem}[Normal Families with Known Variance]
Let $X_1,\ldots,X_n$ be i.i.d. from the $\mathcal{N}(\mu,\sigma^2)$ distribution with an unknown $\mu \in \mathbb{R}$ and a known $\sigma^2$. The joint Lebesgue p.d.f. is
\[f_\mu(x) = \frac{1}{(2\pi\sigma^2)^{n/2}}\exp\left(-\frac{1}{2\sigma^2}\sum(x_i - \mu)^2\right)\]

Then
\[\frac{\partial}{\partial\mu}\log f_\mu(x) = \sum_{i=1}^n (x_i - \mu)/\sigma^2\]

\[I(\mu) = \operatorname{Var}\left(\frac{\sum_{i=1}^n(X_i-\mu)}{\sigma^2}\right) = n\sigma^2/\sigma^4 = n/\sigma^2\]

## Properties of Fisher Information

1. $I(\theta)$ depends on the particular parameterization: If $\theta = \psi(\eta)$ and $\psi$ is differentiable, then the Fisher information about $\eta$ is
\[\tilde{I}(\eta) = \psi'(\eta)^2 I(\psi(\eta))\]
where $I(\theta)$ is the Fisher information about $\theta$.

2. Suppose that $f_\theta$ is twice differentiable in $\theta$ and that
\[\int \frac{\partial^2}{\partial\theta^2}f_\theta(x)I_{f_\theta(x)>0}d\nu = 0\]
Then $I(\theta) = -\mathbb{E}\left[\frac{\partial^2}{\partial\theta^2}\log f_\theta(X)\right]$

3. Regularity condition $(\star)$: $\int \frac{\partial}{\partial\theta}f_\theta(x)d\nu = 0$

4. Let $X$ and $Y$ be independent with the Fisher information about $\theta$ $I_X(\theta)$ and $I_Y(\theta)$, respectively. Suppose $(\star)$ holds for both density functions. Then, the Fisher information about $\theta$ contained in $(X,Y)$ is $I_X(\theta) + I_Y(\theta)$.

5. Fisher information about $\theta$ contained in i.i.d. sample $(X_1,\ldots,X_n)$ is $nI_1(\theta)$ where $I_1(\theta)$ is the Fisher information about $\theta$ contained in a single $X_i$ provided that $(\star)$ holds.
\end{theorem}

\begin{theorem}[Cramér-Rao Lower Bound]
Suppose $\mathcal{P} = \{P_\theta : \theta \in \Theta\}$ satisfies the following conditions:
- $\Theta$ is an open set in $\mathbb{R}$; $P_\theta$ has a p.d.f. $f_\theta$ w.r.t. a measure $\nu$ for all $\theta \in \Theta$
- $f_\theta$ is differentiable as a function of $\theta$ and satisfies
\[0 = \frac{\partial}{\partial\theta} \int f_\theta(x)d\nu = \int \frac{\partial}{\partial\theta} f_\theta(x)d\nu, \quad \theta \in \Theta \quad (1)\]

Suppose that $g(\theta)$ is a differentiable function.
Let $X$ be a sample from $P \in \mathcal{P}$. Suppose $T(X)$ is an unbiased estimator of $g(\theta)$ such that
\[g'(\theta) = \frac{\partial}{\partial\theta} \int T(x) f_\theta(x)d\nu = \int T(x) \frac{\partial}{\partial\theta} f_\theta(x)d\nu, \quad \theta \in \Theta \quad (2)\]

Then $\operatorname{Var}(T(X)) \geq \frac{g'(\theta)^2}{I(\theta)}$ where $I(\theta) > 0$ for any $\theta \in \Theta$.
\end{theorem}

\begin{theorem}[Multivariate C-R Lower Bound]  % Original number: 3.3
When $\theta$ is $k$-dimensional, $g : \Theta \mapsto \mathbb{R}$, the inequality in the Cramér-Rao Lower Bound for $ET = g(\theta)$ $(\theta \in \Theta)$
\[\operatorname{Var}(T(X)) \geq \left[\frac{\partial}{\partial\theta} g(\theta)\right]^T [I(\theta)]^{-1} \frac{\partial}{\partial\theta} g(\theta),\]

where the gradient $\frac{\partial}{\partial\theta} g(\theta) = \left(\frac{\partial}{\partial\theta_1} g(\theta), \ldots, \frac{\partial}{\partial\theta_k} g(\theta)\right)^T$ is assumed to satisfy the regularity condition:
\[\frac{\partial}{\partial\theta} g(\theta) = \int T(x) \frac{\partial}{\partial\theta} f_\theta(x)d\nu.\]

Here we also assume $\mathcal{P} = \{P_\theta : \theta \in \Theta\}$ satisfies the following conditions:
- $\Theta$ is an open set in $\mathbb{R}^k$; $P_\theta$ has a p.d.f. $f_\theta$ w.r.t. a measure $\nu$ for all $\theta \in \Theta$
- $f_\theta$ is differentiable as a function of $\theta$ and satisfies
\[0 = \frac{\partial}{\partial\theta} \int f_\theta(x)d\nu = \int \frac{\partial}{\partial\theta} f_\theta(x)d\nu, \quad \theta \in \Theta.\]
\end{theorem}

\begin{proposition}
Suppose that the distribution of $X$ is from an exponential family $\{f_\theta : \theta \in \Theta\}$, i.e., the p.d.f. of $X$ w.r.t. a $\sigma$-finite measure is
\[f_\theta(x) = \exp\{[\eta(\theta)]^T T(x) - \zeta(\theta)\} h(x),\]
where $\Theta$ is an open subset of $\mathbb{R}^k$.

(i) For any $S$ with $\mathbb{E}|S(X)| < \infty$, it holds that
\[\frac{\partial}{\partial\theta} \int S(x)f_\theta(x)d\nu = \int S(x) \frac{\partial}{\partial\theta} f_\theta(x)d\nu, \quad \theta \in \Theta\]
and
\[I(\theta) = -\mathbb{E}\left[\frac{\partial^2}{\partial\theta\partial\theta^T} \log f_\theta(X)\right].\]
\end{proposition}

\begin{proposition}[Cont.]
(ii) If $I(\eta)$ is the Fisher information matrix for the natural parameter $\eta$, then the covariance matrix $\operatorname{Var}(T) = I(\eta)$.
\end{proposition}

\begin{proposition}[Cont.]
(iii) Let $\psi = \mathbb{E}[T(X)]$. Suppose $I(\psi)$ is the Fisher information matrix for the parameter $\psi$, then $\operatorname{Var}(T) = [I(\psi)]^{-1}$.
\end{proposition}

\end{document}